\clearpage
\setcounter{section}{0}
\setcounter{figure}{0}
\section{Gold implanted in titanium dioxide}
\label{sec:au-tio2-experimental}

\noindent\makebox[\textwidth]{\begin{overpic}[width=\textwidth,height=0.48\textheight,keepaspectratio]{fig-S01.png}
\put(39.922,51.662){\makebox(0,0)[b]{\setlength{\fboxsep}{1pt}\colorbox{black}{\color{white}\fontsize{7}{8}\selectfont 1 nm}}}
\put(89.735,51.662){\makebox(0,0)[b]{\setlength{\fboxsep}{1pt}\colorbox{black}{\color{white}\fontsize{7}{8}\selectfont 1 nm}}}
\put(38.396,3.572){\makebox(0,0)[b]{\setlength{\fboxsep}{1pt}\colorbox{black}{\color{white}\fontsize{7}{8}\selectfont 0.5 nm}}}
\put(88.209,3.572){\makebox(0,0)[b]{\setlength{\fboxsep}{1pt}\colorbox{black}{\color{white}\fontsize{7}{8}\selectfont 0.5 nm}}}
\end{overpic}}
\begingroup\captionsetup{type=figure}
\caption{An experimental image of gold implanted in TiO$_2$ shows improved Blob-Net performance over geometric models, especially at phase boundaries and defects. \textbf{a,} Raw HAADF image of gold (bright region, left) implanted in TiO$_2$ (dark region, right). \textbf{b,} HAADF image processed by mean-field normalization followed by a difference-of-Gaussians approach, used as input to the models. Model outputs are overlaid: Blob-Net (orange), Hex-Net (cyan), and Square-Net (magenta). \textbf{c,} Cropped region visualizing performance at a phase boundary. \textbf{d,} Cropped region visualizing performance around a dislocation core. Further benchmarking on datasets with known ground truth is required to definitively claim better performance on these particular structures, so this example is included in the supplemental information as a visual comparison. Scale bars: a,b, 1 nm; c,d, 0.5 nm.}
\label{fig:au-tio2-three-models}
\endgroup
